\begin{table*}[t]
\caption{Efeito de cada mecanismo no estudo principal: diferenças pareadas entre condições vizinhas}
\label{tab:deltas}
\centering
\footnotesize
\setlength{\tabcolsep}{4pt}
\begin{tabular}{|l|l|c|c|c|}
\hline
\textbf{Comparação} & \textbf{Métrica} & \textbf{GPT-5.6 Luna} & \textbf{Claude Haiku 4.5} & \textbf{Qwen3.5 9B} \\
\hline
RQ1: A $\to$ B & Acurácia factual final & \textbf{$-$0,06 [$-$0,12; $-$0,01]}$^{\dagger}$ & $-$0,04 [$-$0,11; +0,03] & \textbf{$-$0,07 [$-$0,13; $-$0,02]}$^{\dagger}$ \\
 & Retenção & \textbf{$-$0,09 [$-$0,17; $-$0,02]}$^{\dagger}$ & $-$0,06 [$-$0,17; +0,04] & +0,01 [$-$0,07; +0,10] \\
 & Taxa de correção$^{a}$ & +0,04 [$-$0,03; +0,10] & +0,03 [$-$0,06; +0,11] & +0,01 [$-$0,07; +0,09] \\
 & ESR & $-$0,03 [$-$0,09; +0,02] & $-$0,05 [$-$0,10; +0,01] & $-$0,02 [$-$0,08; +0,04] \\
 & Profundidade & \textbf{+0,43 [+0,32; +0,54]} & \textbf{+0,26 [+0,15; +0,37]} & \textbf{+0,31 [+0,22; +0,41]} \\
 & Amplificação & $-$0,17 [$-$0,42; +0,07] & \textbf{$-$0,66 [$-$0,90; $-$0,44]} & $-$0,35 [$-$0,75; +0,04] \\
 & RCCA & +0,01 [$-$0,05; +0,07] & $-$0,02 [$-$0,09; +0,04] & \textbf{+0,22 [+0,16; +0,28]} \\
 & FESA líquida$^{b}$ & $-$0,03 [$-$0,10; +0,04] & \textbf{$-$0,09 [$-$0,16; $-$0,02]}$^{\dagger}$ & \textbf{+0,16 [+0,05; +0,28]}$^{\dagger}$ \\
\hline
RQ2: B $\to$ C & Acurácia factual final & \textbf{+0,11 [+0,05; +0,16]} & \textbf{+0,11 [+0,05; +0,18]} & $-$0,03 [$-$0,07; +0,02] \\
 & Retenção & \textbf{+0,17 [+0,09; +0,24]} & \textbf{+0,19 [+0,09; +0,29]} & +0,00 [$-$0,06; +0,06] \\
 & Taxa de correção$^{a}$ & \textbf{+0,15 [+0,07; +0,22]} & \textbf{+0,18 [+0,09; +0,28]} & +0,00 [$-$0,06; +0,06] \\
 & ESR & +0,03 [$-$0,03; +0,09] & \textbf{+0,07 [+0,02; +0,13]}$^{\dagger}$ & +0,01 [$-$0,05; +0,08] \\
 & Profundidade & \textbf{+0,12 [+0,07; +0,18]} & \textbf{+0,17 [+0,12; +0,24]} & +0,05 [+0,00; +0,11] \\
 & Amplificação & +0,18 [$-$0,05; +0,40] & \textbf{+0,60 [+0,36; +0,86]} & +0,32 [$-$0,01; +0,66] \\
 & RCCA & \textbf{+0,15 [+0,09; +0,20]} & \textbf{+0,07 [+0,02; +0,13]}$^{\dagger}$ & \textbf{+0,08 [+0,02; +0,15]}$^{\dagger}$ \\
 & FESA líquida$^{b}$ & \textbf{+0,16 [+0,09; +0,23]} & \textbf{+0,09 [+0,01; +0,17]}$^{\dagger}$ & $-$0,03 [$-$0,12; +0,07] \\
\hline
RQ3: C $\to$ D & Acurácia factual final & \textbf{+0,15 [+0,08; +0,21]} & \textbf{+0,10 [+0,04; +0,16]} & \textbf{+0,10 [+0,05; +0,14]} \\
 & Retenção & \textbf{+0,17 [+0,09; +0,25]} & \textbf{+0,14 [+0,05; +0,24]} & \textbf{+0,11 [+0,01; +0,21]}$^{\dagger}$ \\
 & Taxa de correção$^{a}$ & \textbf{+0,16 [+0,08; +0,24]} & +0,09 [+0,00; +0,17] & \textbf{+0,10 [+0,02; +0,18]}$^{\dagger}$ \\
 & ESR & \textbf{$-$0,10 [$-$0,16; $-$0,04]} & \textbf{$-$0,14 [$-$0,21; $-$0,08]} & $-$0,05 [$-$0,11; +0,02] \\
 & Profundidade & \textbf{$-$0,12 [$-$0,19; $-$0,05]} & \textbf{$-$0,16 [$-$0,23; $-$0,09]} & $-$0,05 [$-$0,11; +0,01] \\
 & Amplificação & $-$0,04 [$-$0,27; +0,21] & \textbf{$-$0,41 [$-$0,65; $-$0,18]} & +0,20 [$-$0,28; +0,86] \\
 & RCCA & \textbf{+0,43 [+0,37; +0,49]} & \textbf{+0,35 [+0,29; +0,41]} & +0,01 [$-$0,05; +0,07] \\
 & FESA líquida$^{b}$ & \textbf{+0,45 [+0,38; +0,52]} & \textbf{+0,31 [+0,24; +0,38]} & \textbf{+0,10 [+0,01; +0,20]}$^{\dagger}$ \\
\hline
\end{tabular}
\\[2pt]{\scriptsize\raggedright Diferença média pareada por pergunta (segunda condição menos a primeira) e IC95\% bootstrap. Em negrito: intervalo que não inclui zero; $^{\dagger}$: não sobrevive à correção de Holm sobre as 24 comparações do modelo. Para ESR e profundidade, descontado o piso de ruído do teste A/A (Tabela~\ref{tab:verificacoes}), só o aumento de profundidade de A para B no GPT e no Haiku continua fora de zero. $^{a}$Calculada só sobre as perguntas elegíveis nas duas condições comparadas, por isso não é a diferença das médias da Tabela~\ref{tab:principal}. $^{b}$FESA menos o placebo da mesma pergunta e condição.\par}
\end{table*}
